\documentclass[11pt]{article}
\usepackage{../shared/luxcover}
\usepackage{amsmath,amssymb}
\usepackage{hyperref}
\usepackage{booktabs}
\usepackage{enumitem}
\usepackage{xcolor}

\title{Lux Z-Chain: Zero-Knowledge Proof Chain}
\author{Zach Kelling\\\textit{Lux Industries, Inc.}\\\texttt{research@lux.network}}
\date{May 2023}

\begin{document}

\luxcoverpage

\begin{abstract}
Quantum-Resilient Privacy and Scalability
\end{abstract}

\tableofcontents
\newpage


\section{Introduction}

ZChain is a revolutionary blockchain platform that prioritizes privacy, scalability, and quantum resistance. This white paper provides an overview of ZChain's core concepts, architecture, and unique features that set it apart in the blockchain ecosystem.


Zero-knowledge proof systems, such as ZChain, offer a compelling solution for private and scalable blockchain applications. ZChain's efficient design allows it to handle a higher volume of transactions compared to traditional proof systems while ensuring privacy remains intact. Moreover, ZChain's compatibility with Ethereum opens the door to building private and scalable versions of existing Ethereum applications.


Private and scalable blockchain applications are crucial for businesses as they enable safeguarding of sensitive data, adherence to regulatory requirements like GDPR, and reduction of infrastructure costs. With data protection becoming increasingly important for businesses, private and scalable blockchain applications emerge as a robust solution.


\section{Zero Knowledge Computing: A Primer}

Zero-Knowledge (ZK) computing is a powerful cryptographic technique that allows a prover to demonstrate the truth of a statement to a verifier without divulging any additional information beyond the statement's validity. In simpler terms, ZK computing enables the prover to convince the verifier that a statement is true without revealing the underlying data or computation involved.


At the core of ZK computing are ZK circuits, fundamental building blocks that represent computations or a series of computations. These circuits consist of logical gates, such as AND, OR, and NOT gates, which manipulate input variables to perform various operations. ZK circuits can be designed to handle a wide range of computations, from basic arithmetic to intricate algorithms.


The strength of ZK circuits lies in their ability to process and validate data while preserving privacy. By utilizing ZK circuits, it becomes possible to verify the accuracy of computations without exposing the actual data being processed. This characteristic makes ZK circuits an essential tool for applications in blockchain technology that prioritize privacy.


\section{zkEVM and Privacy-Enhanced Smart Contracts}

Through zkEVM, developers can construct complex computations while keeping the input data concealed. This enables secure execution of smart contracts involving private or sensitive information, such as financial transactions or personal data. By leveraging ZK circuits and zero-knowledge proofs, zkEVM empowers developers to build decentralized applications (DApps) with advanced privacy features.


How it works:


A zkEVM replicates the Ethereum environment as a zero-knowledge rollup, bringing the Ethereum developer experience and existing tooling to a highly scalable and secure layer 2.


The EVM executes smart contracts and computes the state of the Lux network after each new block is added to the chain.


The zkEVM takes the initial state, computes transactions, and outputs a new, updated state and an accompanying zero-knowledge proof.


zkEVMs leverage zero-knowledge encryption and proofs to allow smart contracts to run with minimal data, and it can encrypt sensitive information in cryptocurrency transactions.


\section{Scalability through Rollups}

ZChain achieves scalability by leveraging zero-knowledge rollups, a mechanism that separates transaction execution from data and performs complex calculations off-chain. In a zero-knowledge rollup, transactions are bundled together and summarized into a single proof that can be easily validated by third-party observers. This approach significantly reduces the number of transactions that need to be verified on the main blockchain, such as Lux, resulting in improved computational efficiency and scalability.


ZChain's zero-knowledge rollups enable the network to handle large volumes of transactions while preserving privacy and efficiency.


\section{Privacy Focus}

One of the key features of ZChain is its ability to ensure privacy in transactions. Through the use of advanced cryptographic techniques, ZChain enables users to conceal the identity and content of their transactions. This is achieved through a combination of zero-knowledge proofs and blind signatures. Blind signatures, a technique pioneered by David Chaum, allow users to obtain a signature on a message without revealing the content of the message to the signer. Zero-knowledge proofs, on the other hand, allow users to prove the validity of a statement without revealing any additional information. By leveraging these techniques, ZChain ensures that transactions remain private and confidential.


\section{Compliance with Regulations}

In an era of stringent data protection regulations, like the General Data Protection Regulation (GDPR), private and scalable blockchain applications become essential tools for compliance. By employing ZChain's privacy-enhancing capabilities, businesses can demonstrate adherence to regulatory requirements by ensuring that only authorized users have access to personal and sensitive data. With ZChain, businesses can establish auditable and immutable records of data access, offering transparency and accountability that aligns with regulatory expectations. This compliance-friendly approach empowers businesses to confidently navigate complex regulatory landscapes while protecting the privacy of their users' data.


\section{Cost Reduction and Efficiency}

Private and scalable blockchain applications present an opportunity for businesses to achieve cost savings and operational efficiency. By leveraging ZChain, businesses can reduce reliance on expensive infrastructure, such as data centers, by utilizing the decentralized nature of the blockchain. Private blockchains offer a secure and cost-effective alternative, eliminating the need for physical infrastructure and reducing associated maintenance costs. Additionally, the scalability of ZChain allows businesses to handle a higher transaction volume without compromising privacy, resulting in streamlined operations and improved cost efficiency. This cost-saving potential extends to security measures as well, as private and scalable blockchain applications alleviate the need for physical security measures like security guards and cameras.


\section{Quantum Safe}

As quantum computing technology progresses, conventional cryptographic algorithms employed in blockchain networks become susceptible to attacks. In response to this challenge, ZChain has implemented quantum-safe security measures.


A crucial aspect of ZChain's quantum resistance strategy revolves around the adoption of lattice-based encryption. Lattice-based cryptography provides robust protection against quantum computer attacks, thanks to its resilient mathematical foundation. By employing lattice encryption at both the account and consensus levels, ZChain guarantees the security of user data and the overall network, safeguarding them against quantum threats.


\section{Empowering Web3}

ZK-scaling on Lux will enable improved functionality and UX for virtually every web 3 application, with particular benefits for DAO tooling, DEXs, and NFT management.


By leveraging zero-knowledge proofs, Lux overcomes the scalability limitations of traditional blockchain networks. This enables seamless and efficient execution of complex operations, such as DAO governance, trading on DEXs, and managing Non-Fungible Tokens (NFTs). With Lux, users can expect reduced costs, faster transaction speeds, and an overall enhanced experience when interacting with web3 applications.


Lux's native support for NFTs opens up new possibilities for digital asset management. Users can effortlessly buy, sell, and trade any type of digital asset on the Ethereum network. Lux's seamless integration with Ethereum enables the creation of fully composable NFTs, unlocking innovative use cases such as DAO membership tracking and reputation building through blockchain-based credit scores. With Lux, NFTs become more than just digital collectibles; they become powerful tools for decentralized governance and identity verification.


DAOs play a pivotal role in the Web3 landscape, enabling decentralized decision-making and community governance. Lux's scalability and privacy features provide a solid foundation for DAO tooling, empowering communities to efficiently manage voting, treasury, and governance processes. Additionally, Lux's improved transaction throughput and low fees make it an ideal platform for decentralized exchanges, enhancing liquidity and enabling seamless trading experiences for users.


\section{Private Applications}

Zero-knowledge proof systems like ZChain are particularly well-suited for private and scalable blockchain applications. ZChain's efficient design means that it can handle more transactions than traditional proof systems, without sacrificing privacy. Additionally, ZChain's compatibility with Ethereum means that it can be used to build private and scalable versions of existing Ethereum applications. Private and scalable blockchain applications are important for a number of reasons. First, they can help businesses protect sensitive data. Second, they can help businesses comply with regulations, like GDPR. Finally, they can help businesses save costs by reducing the need for expensive infrastructure. Businesses are increasingly looking for ways to protect their data. Private and scalable blockchain applications can help businesses keep their data safe by ensuring that only authorized users can access it. Additionally, private blockchains can help businesses comply with regulations, like GDPR, by ensuring that only authorized users can access sensitive data. Private and scalable blockchain applications can also help businesses save costs.


\subsection{Unified Programming Environment}

Lux's cutting-edge technology and unified development ecosystem, coupled with efficient ZK-scaling, are set to revolutionize blockchain programming. Developers will now have the unprecedented ability to construct complex applications by assembling individual smart contracts that have been ZK-verified by a single trusted prover, all while enjoying the advantages of a functional, typed programming language.


This groundbreaking approach offers a multitude of benefits, including speed, flexibility, and enhanced security. It paves the way for a new era in blockchain development, empowering developers to create previously unimaginable applications. With Lux's advanced capabilities, complex gaming physics simulations, virtual reality engines, expansive artistic sandboxes, and a wide range of other innovative applications become a reality.


By leveraging Lux's composable and unified development ecosystem, developers can seamlessly combine verified smart contracts to construct sophisticated and interconnected systems. The ZK-verification process ensures the integrity and correctness of each contract, providing a high level of security and trustworthiness.


The integration of a functional, typed language further enhances the development experience. Developers can leverage the benefits of strong typing and functional programming paradigms, resulting in cleaner and more reliable code. This facilitates easier debugging, maintainability, and code reuse, ultimately accelerating the development process and reducing the likelihood of errors.


The speed and flexibility offered by Lux's technology stack enable developers to push the boundaries of what is possible in blockchain programming. The elimination of scalability limitations allows for the creation of resource-intensive applications that demand high computational power and extensive data processing. This opens up a vast array of possibilities in areas such as gaming, virtual reality, digital art, and beyond.


Lux's vision is to empower developers to unlock the full potential of blockchain technology. By providing the necessary tools, scalability, and security, Lux is spearheading a new frontier in blockchain development. The era of complex, cutting-edge applications has arrived, and Lux is at the forefront, leading the way towards a future of limitless possibilities.


\section{Ethereum Bootstrapping}

ZChain, as a ZK-rollup solution, has the capability to settle transactions on any blockchain. However, for practical purposes, settlement requires a robust and decentralized network with full programmability, making Ethereum an ideal choice. Ethereum's settlement layer serves as the ultimate source of truth, and its infrastructure's security and trustworthiness play a crucial role in the execution and data availability mechanisms.


By leveraging the security guarantees and extensive network of web 3 resources offered by Ethereum, ZChain operates as a ZK-rollup to Ethereum. This integration allows Lux to capitalize on Ethereum's established ecosystem while benefiting from its vast network effects. Lux has already developed its own Layer 1 blockchain and data availability solutions, enabling seamless transactions between all supported networks.


\section{Conclusion}

ZChain represents a significant advancement in the blockchain industry, offering a platform that combines privacy, scalability, and quantum resistance. By leveraging advanced cryptographic techniques such as zero-knowledge proofs, blind signatures, and lattice encryption, ZChain provides users with a secure and private environment for conducting transactions.


The innovative approach of ZChain sets it apart from traditional blockchain platforms and paves the way for a new era of privacy-preserving and quantum-safe blockchain technology.


\section{References}

D. Chaum, "Blind signatures for untraceable payments," Advances in Cryptology - CRYPTO '82, pp. 199-203, 1983.


S. Goldwasser et al., "Probabilistic encryption \& how to play mental poker keeping secret all partial information," Advances in Cryptology - CRYPTO '82, pp. 465-491, 1983.


C. Peikert, "A decade of lattice cryptography," Foundations and Trends in Theoretical Computer Science, vol. 10, no. 4, pp. 283-424, 2016.


M. Ben-Or et al., "The Byzantine Generals Problem," ACM Transactions on Programming Languages and Systems, vol. 4, no. 3, pp. 382-401, 1982.


G. Wood, "Ethereum: A secure decentralised generalised transaction ledger," Ethereum Project Yellow Paper, vol. 151, 2014.


D. Mazières et al., "The Stellar Consensus Protocol: A Federated Model for Internet-level Consensus," Stellar Development Foundation, 2015.


E. Ghassan et al., "HotStuff: BFT Consensus in the Lens of Blockchain," Proceedings of the 39th Symposium on Principles of Distributed Computing, 2020.


E. Syta et al., "Keeping Authorities 'Honest or Bust' with Decentralized Witness Cosigning," Proceedings of the 2016 ACM SIGSAC Conference on Computer and Communications Security, pp. 526-537, 2016.


A. E. Kosba et al., "Hawk: The Blockchain Model of Cryptography and Privacy-Preserving Smart Contracts," Proceedings of the 2016 ACM SIGSAC Conference on Computer and Communications Security, pp. 839-851, 2016.


B. Bünz et al., "Transparent SNARKs from DARK Compilers," Advances in Cryptology - EUROCRYPT '20, pp. 1-32, 2020.


V. Buterin, "Ethereum: A Next-Generation Smart Contract and Decentralized Application Platform," White Paper, 2013.


G. Wood, "Ethereum: A Secure Decentralized Generalized Transaction Ledger," Yellow Paper, 2014.


R. Meher, A. Sinha, and S. Ghosh, "Ethereum Virtual Machine: A Blockchain-Based Distributed Computing Platform," Proceedings of the 2019 2nd International Conference on Intelligent Sustainable Systems (ICISS), pp. 498-503, 2019.


L. Luu, V. Narayanan, C. Zheng, K. Baweja, S. Gilbert, and P. Saxena, "A Secure Sharding Protocol for Open Blockchains," Proceedings of the 2016 ACM SIGSAC Conference on Computer and Communications Security, pp. 17-30, 2016.


\end{document}
